\documentclass[12pt,a4paper]{article}
\usepackage[margin=2.5cm]{geometry}
\usepackage{amsmath,amssymb,amsthm}
\usepackage{physics}
\usepackage{booktabs}
\usepackage{hyperref}
\usepackage{xcolor}
\usepackage{float}

\hypersetup{colorlinks=true, linkcolor=blue, citecolor=blue, urlcolor=blue}

\newcommand{\psieff}{\psi_{\mathrm{eff}}}
\newcommand{\neff}{n_{\mathrm{eff}}}
\newcommand{\etac}{\eta_{\mathrm{c}}}

\title{DFD Propulsion Forces:\\
The Correct Geodesic Interpretation\\[6pt]
\large All Forces Through the Optical Metric, No $G/c^4$ Suppression}
\author{DFD Research Programme}
\date{March 2026}

\begin{document}
\maketitle

\begin{abstract}
Previous analyses of DFD-based propulsion incorrectly separated ``optical'' from ``mechanical'' effects.
In Density Field Dynamics, there \emph{is} no separation: all matter follows geodesics of the optical
metric $\tilde{g}_{\mu\nu}$, defined by the refractive index $n = e^\psi$. The above-threshold
modification $\neff = \exp[\psi + \kappa(\eta - \etac)\,\Theta(\eta - \etac)]$ enters the geodesic
equation \emph{directly}, producing acceleration $a = (c^2/2)\,\nabla(\ln \neff)$ with \textbf{no}
$G/c^4$ suppression factor. We compute forces for six scenarios ranging from pulsed laboratory magnets
to focused petawatt lasers. The raw numbers are enormous---indicating either a transformative
experimental signature or the onset of nonlinear saturation that must be resolved theoretically.
We analyse the saturation question and the self-force problem in detail.
\end{abstract}

\tableofcontents
\newpage

%=======================================================================
\section{The Founding Postulate: Geodesic Motion in the Optical Metric}
\label{sec:postulate}
%=======================================================================

The DFD action for a point particle is (cf.\ \texttt{section\_formalism.tex}, line~178):
\begin{equation}
\label{eq:Spp}
S_{\mathrm{pp}} = -mc \int d\tau \sqrt{-\tilde{g}_{\mu\nu}\,\frac{dx^\mu}{d\tau}\,\frac{dx^\nu}{d\tau}}\,,
\end{equation}
where the optical metric is
\begin{equation}
\tilde{g}_{\mu\nu} = n^2\, g_{\mu\nu}\,,\qquad n = e^\psi\,.
\end{equation}
Matter follows geodesics of $\tilde{g}_{\mu\nu}$. In the weak-field, non-relativistic limit,
the geodesic equation gives the acceleration:
\begin{equation}
\label{eq:accel}
\boxed{\vec{a} = \frac{c^2}{2}\,\nabla(\ln n^2) = c^2\,\nabla\psi\,.}
\end{equation}
More precisely, using $\ln n = \psi$, the acceleration magnitude is
\begin{equation}
a = \frac{c^2}{2}\,|\nabla(\ln n^2)| = c^2\,|\nabla\psi|\,.
\end{equation}
(The factor depends on the convention for $n$ vs.\ $n^2$ in the conformal factor. We use
$a = (c^2/2)\,|\nabla(\ln \neff^2)| = c^2\,|\nabla\psieff|$ throughout, which for the threshold
contribution gives $a = (c^2/2)\,|\nabla[\ln \neff^2]|$ where $\ln \neff = \psieff$.)

\medskip\noindent
\textbf{Critical point:} This acceleration is \emph{not} derived from the gravitational field
equation. It comes from the \emph{geodesic equation}. The field equation (with its $G/c^4$ coupling)
determines $\psi$ from sources; but once $\psi$ is known---by any mechanism---the acceleration
follows from Eq.~\eqref{eq:accel} with \emph{no} additional suppression factor.

%=======================================================================
\section{The Above-Threshold Mechanism}
\label{sec:threshold}
%=======================================================================

When the electromagnetic energy ratio
\begin{equation}
\eta \;\equiv\; \frac{U_{\mathrm{EM}}}{\rho_m\, c^2}
\end{equation}
exceeds the threshold $\etac = \alpha/4 = 1.824 \times 10^{-3}$, the effective refractive index
acquires an additional contribution:
\begin{equation}
\label{eq:neff}
\neff = \exp\!\bigl[\psi + \kappa\,(\eta - \etac)\,\Theta(\eta - \etac)\bigr]\,,
\end{equation}
where
\begin{equation}
\kappa = k_a = \frac{3}{8\alpha} = 51.39\,,\qquad
\etac = \frac{\alpha}{4} = 1.824 \times 10^{-3}\,.
\end{equation}
The above-threshold modification of the log-refractive index is
\begin{equation}
\delta\psieff = \kappa\,(\eta - \etac)\,.
\end{equation}
The acceleration from a gradient in $\delta\psieff$ is
\begin{equation}
\label{eq:a_threshold}
\boxed{a = \frac{c^2}{2}\,|\nabla(\delta\psieff \cdot 2)| = c^2\,|\nabla[\kappa(\eta - \etac)]|\,.}
\end{equation}

\noindent\textbf{There is no $G/c^4$ suppression.} The coupling constant is $\kappa \approx 51.4$,
which is of order unity, not of order $G/c^4 \sim 10^{-44}$. The $c^2$ factor in the geodesic
equation is \emph{enhancing}, not suppressing.

%=======================================================================
\section{Fundamental Parameters}
\label{sec:params}
%=======================================================================

\begin{table}[H]
\centering
\begin{tabular}{lll}
\toprule
\textbf{Parameter} & \textbf{Symbol} & \textbf{Value} \\
\midrule
Fine structure constant & $\alpha$ & $7.297 \times 10^{-3}$ \\
Threshold coupling & $\kappa = 3/(8\alpha)$ & $51.39$ \\
Critical EM ratio & $\etac = \alpha/4$ & $1.824 \times 10^{-3}$ \\
Speed of light & $c$ & $2.998 \times 10^8$ m/s \\
Vacuum permeability & $\mu_0$ & $1.257 \times 10^{-6}$ T$\cdot$m/A \\
\bottomrule
\end{tabular}
\caption{DFD parameters used throughout.}
\end{table}

For a magnetic field $B$, the EM energy density is $U_{\mathrm{EM}} = B^2/(2\mu_0)$, and
\begin{equation}
\eta = \frac{B^2}{2\mu_0\,\rho\, c^2}\,.
\end{equation}
The threshold condition $\eta > \etac$ requires
\begin{equation}
\label{eq:rho_threshold}
\rho < \rho_{\mathrm{th}} = \frac{B^2}{2\mu_0\,\etac\, c^2}\,.
\end{equation}

%=======================================================================
\section{Scenario 1: Pulsed Magnet in Moderate Vacuum}
\label{sec:scenario1}
%=======================================================================

\subsection{Setup}

\begin{itemize}
\item $B = 50$ T (pulsed magnet, NHMFL Los Alamos class)
\item $\rho = 10^{-7}$ kg/m$^3$ (rough vacuum, easily achievable)
\item Bore radius $R = 0.05$ m, length $L = 0.2$ m
\end{itemize}

\subsection{Computation}

\begin{align}
U_{\mathrm{EM}} &= \frac{B^2}{2\mu_0} = \frac{2500}{2 \times 1.257 \times 10^{-6}} = 9.947 \times 10^8 \;\text{J/m}^3 \\[6pt]
\eta &= \frac{9.947 \times 10^8}{10^{-7} \times (2.998 \times 10^8)^2} = 0.1107 \\[6pt]
\frac{\eta}{\etac} &= \frac{0.1107}{1.824 \times 10^{-3}} = 60.7 \\[6pt]
\delta\psieff &= \kappa\,(\eta - \etac) = 51.39 \times 0.1089 = 5.594 \\[6pt]
\exp(\delta\psieff) &= 268.7
\end{align}

The refractive index changes by a factor of $\sim 270$ inside the field region relative to outside.

\subsection{Gradient and force at the bore edge}

Inside a uniform solenoid, $B$ is approximately uniform, so $\eta$ is uniform (assuming uniform $\rho$).
Therefore:
\begin{equation}
\nabla\eta = 0 \quad\text{inside the bore}\,.
\end{equation}
\textbf{There is no force on a test mass in the interior of a uniform-field region.}
The gradient exists \emph{only} at the boundary, where the field drops from $B$ to 0 over the
fringe-field length, which is approximately the bore radius $R$.

At the bore edge:
\begin{align}
|\nabla(\delta\psieff)| &\approx \frac{\delta\psieff}{R} = \frac{5.594}{0.05} = 111.9 \;\text{m}^{-1} \\[6pt]
a &= \frac{c^2}{2} \times 111.9 = 5.03 \times 10^{18} \;\text{m/s}^2 \\[6pt]
F_{\text{1\,kg}} &= 5.03 \times 10^{18} \;\text{N}
\end{align}

\subsection{Assessment}

The raw number---$5 \times 10^{18}$ m/s$^2$, or $5 \times 10^{17}\,g$---is \emph{absurdly} large.
This is $10^{17}$ times stronger than any force ever measured in a laboratory.
Several interpretations:

\begin{enumerate}
\item The linear threshold formula $\delta\psieff = \kappa(\eta - \etac)$ breaks down at
      $\eta/\etac = 60$ (see Section~\ref{sec:saturation}).
\item The ``force'' applies to the fringe region only, and the actual measurable effect on a
      test mass depends on integration over the trajectory through the fringe.
\item If even a fraction of this force is real, it would be \emph{instantly} detectable---and
      would likely have been noticed already in high-field pulsed-magnet experiments, unless
      the vacuum condition is critical.
\item The vacuum condition $\rho = 10^{-7}$ kg/m$^3$ is key: standard pulsed magnets operate
      in air ($\rho \sim 1.2$ kg/m$^3$), giving $\eta/\etac \sim 5 \times 10^{-9}$, far below
      threshold.
\end{enumerate}

%=======================================================================
\section{Scenario 2: Asymmetric Magnet Configuration}
\label{sec:scenario2}
%=======================================================================

\subsection{Setup}

An asymmetric configuration with $B_{\text{high}} = 50$ T on one side and $B_{\text{low}} = 0$ T on
the other, separated by $d = 0.3$ m, in $\rho = 10^{-7}$ kg/m$^3$ vacuum.

\subsection{Computation}

\begin{align}
\delta\psieff^{\text{high}} &= 5.594 \quad (\text{from Scenario 1}) \\
\delta\psieff^{\text{low}} &= 0 \quad (\text{below threshold}) \\
\nabla\psieff &\approx \frac{5.594 - 0}{0.3} = 18.6 \;\text{m}^{-1} \\
a &= \frac{c^2}{2} \times 18.6 = 8.38 \times 10^{17} \;\text{m/s}^2
\end{align}

For a 100~kg device entirely within the gradient: $F = 8.4 \times 10^{19}$~N.

\subsection{The self-force question}
\label{sec:selfforce_scenario2}

The numbers above assume that the device experiences its own $\psieff$ gradient as an external force.
This requires careful analysis---see Section~\ref{sec:selfforce}.

%=======================================================================
\section{Scenario 3: Focused Laser in High Vacuum}
\label{sec:scenario3}
%=======================================================================

\subsection{Setup}

\begin{itemize}
\item $I = 10^{18}$ W/m$^2$ (petawatt laser, e.g.\ ELI, OMEGA EP)
\item Beam waist $w_0 = 5\;\mu$m
\item $\rho = 10^{-10}$ kg/m$^3$ (ultra-high vacuum)
\end{itemize}

\subsection{Computation}

For a plane wave, $U_{\mathrm{EM}} = I/c$:
\begin{align}
U_{\mathrm{EM}} &= \frac{10^{18}}{2.998 \times 10^8} = 3.336 \times 10^9 \;\text{J/m}^3 \\[6pt]
\eta &= \frac{3.336 \times 10^9}{10^{-10} \times (2.998 \times 10^8)^2} = 371.1 \\[6pt]
\frac{\eta}{\etac} &= 2.03 \times 10^5 \\[6pt]
\delta\psieff &= 51.39 \times 371.1 = 19{,}071
\end{align}

This is \emph{vastly} beyond any regime where the linear threshold formula could hold.
The formal gradient gives:
\begin{align}
|\nabla(\delta\psieff)| &\approx \frac{19{,}071}{5 \times 10^{-6}} = 3.81 \times 10^9 \;\text{m}^{-1} \\[6pt]
a &= \frac{c^2}{2} \times 3.81 \times 10^9 = 1.71 \times 10^{26} \;\text{m/s}^2
\end{align}

\subsection{Assessment}

These numbers are unphysical in the linear approximation. A $\delta\psieff$ of 19,071 corresponds to
$\exp(\delta\psieff) \sim 10^{8{,}283}$---a refractive index change that is nonsensical. The linear
threshold formula \emph{certainly} breaks down here. The laser scenario is deep into the nonlinear/saturation
regime and cannot be analysed without the full nonlinear theory.

However, the laser scenario demonstrates that the \emph{threshold condition} $\eta > \etac$ is easily
satisfied with focused high-intensity lasers in vacuum. Even at much lower intensities, $\eta/\etac \gg 1$.
The question is what happens at these extreme ratios.

%=======================================================================
\section{Scenario 4: Saturation Analysis}
\label{sec:saturation}
%=======================================================================

\subsection{The linear regime and its limits}

The formula $\delta\psieff = \kappa(\eta - \etac)$ is a linear approximation valid near threshold.
Table~\ref{tab:saturation} shows how $\delta\psieff$ and $\exp(\delta\psieff)$ grow with $\eta/\etac$:

\begin{table}[H]
\centering
\begin{tabular}{rrrr}
\toprule
$\eta/\etac$ & $\delta\psieff$ & $\exp(\delta\psieff)$ & Regime \\
\midrule
2   & 0.094  & 1.098     & Linear, small perturbation \\
5   & 0.375  & 1.455     & Linear, moderate \\
10  & 0.844  & 2.325     & \textbf{Solar corona calibration} \\
20  & 1.781  & 5.937     & Mildly nonlinear \\
60  & 5.531  & 252.5     & \textbf{50~T magnet scenario} \\
100 & 9.281  & $1.07 \times 10^4$ & Strongly nonlinear \\
370 & 34.59  & $1.06 \times 10^{15}$ & Extremely nonlinear \\
\bottomrule
\end{tabular}
\caption{Growth of $\delta\psieff$ with $\eta/\etac$. The solar corona light-bending calibration
operates at $\eta/\etac \sim 10$.}
\label{tab:saturation}
\end{table}

\subsection{The calibration anchor}

The DFD framework is calibrated against the solar corona, where $\eta/\etac \sim 10$ and $\kappa = 51.4$
yields $\delta\psieff \approx 0.84$, giving $\exp(\delta\psieff) \approx 2.3$. This is the regime where
the theory is \emph{known} to work. At this point, the refractive index modification is a factor of 2---large,
but not unphysical.

\subsection{Beyond the calibrated regime}

At $\eta/\etac = 60$ (the 50~T magnet scenario), $\delta\psieff = 5.6$ and the refractive index changes
by a factor of 270. This is 6.6 times the calibrated $\delta\psieff$ value. Several possibilities:

\begin{enumerate}
\item \textbf{No saturation:} The linear formula holds, and the enormous forces are real. This would
      mean DFD predicts qualitatively new physics at high $\eta$, testable with pulsed magnets in vacuum.

\item \textbf{Logarithmic saturation:} $\delta\psieff$ saturates as $\sim \kappa\,\etac\ln(\eta/\etac)$
      at large $\eta$. At $\eta/\etac = 60$: $\delta\psieff \sim 51.4 \times 1.82 \times 10^{-3} \times
      \ln(60) = 0.094 \times 4.09 = 0.38$, giving $\exp(\delta\psieff) = 1.47$. Forces would be
      reduced by a factor of $\sim 15$ from the linear prediction.

\item \textbf{Hard saturation at $\delta\psieff^{\max}$:} The modification saturates at some maximum
      value, perhaps $\delta\psieff^{\max} \sim 1$. Forces are bounded.

\item \textbf{Nonlinear backreaction:} At large $\delta\psieff$, the modified metric feeds back into
      the field equation, limiting further growth. This is analogous to how large gravitational potentials
      in GR resist further deepening.
\end{enumerate}

\subsection{What the solar corona tells us}

At $\eta/\etac = 10$, the linear formula gives $\delta\psieff = 0.84$, which \emph{successfully}
predicts the observed light deflection. This means:
\begin{itemize}
\item The formula works at $\delta\psieff \sim 1$.
\item We have no direct evidence for $\delta\psieff \gg 1$.
\item The 50~T scenario ($\delta\psieff = 5.6$) is extrapolation, not interpolation.
\item Conservative approach: trust the formula up to $\delta\psieff \sim 1$--$2$, and flag
      larger values as requiring nonlinear analysis.
\end{itemize}

\subsection{Practical implication}

Even with saturation, the forces near threshold ($\eta/\etac \sim 2$--$5$) give:
\begin{equation}
\delta\psieff \sim 0.1\text{--}0.4\,,\quad a \sim \frac{c^2}{2}\frac{\delta\psieff}{R}
\sim 10^{15}\text{--}10^{16}\;\text{m/s}^2\,.
\end{equation}
This is still enormous. The $c^2/2$ prefactor ($4.5 \times 10^{16}$ m$^2$/s$^2$) ensures that
\emph{any} non-negligible $\delta\psieff$ gradient over laboratory scales produces colossal accelerations.

\textbf{If DFD is correct and the threshold mechanism operates as described, the effect is either
(a) absurdly large and immediately detectable, or (b) saturated/suppressed by physics not yet
captured in the linear formula.}

%=======================================================================
\section{Scenario 5: The Self-Force Problem}
\label{sec:selfforce}
%=======================================================================

\subsection{Statement of the problem}

For propulsion, the device must accelerate itself using its own EM field. The EM field creates a
gradient in $\neff$, and the question is whether the device accelerates in this self-generated gradient.

\subsection{Scenario (A): Universal coupling (equivalence principle)}

DFD posits that \emph{all} matter follows geodesics of $\tilde{g}_{\mu\nu}$. If the device creates
a $\psieff$ gradient, then:
\begin{itemize}
\item All matter in the gradient region accelerates equally (WEP).
\item The device, its occupants, and any test masses all accelerate together.
\item There are no g-forces on occupants---this is the ``$\psi$-bubble'' concept.
\end{itemize}

This is analogous to gravitational freefall. However, a critical difference:

\textbf{In gravity:} A mass does not accelerate in its own field. The centre of mass of an isolated
system does not move due to internal gravitational forces. This is guaranteed by Newton's third law
and momentum conservation.

\textbf{In DFD:} The $\psieff$ modification is \emph{not} sourced through the gravitational field equation
(which carries the $G/c^4$ suppression). It is a \emph{direct} modification of the optical metric by EM
energy above threshold. The question is whether the analogous self-force cancellation still applies.

\subsection{Momentum conservation analysis}

Total momentum of the system:
\begin{equation}
P_{\mathrm{total}} = P_{\mathrm{matter}} + P_{\mathrm{EM}} + P_\psi\,.
\end{equation}

If total momentum is conserved (as it must be in any consistent theory), then for the centre of mass
to accelerate, something must carry the recoil momentum.

\subsubsection{The $\psi$-field momentum}

Agent~6 derived the $\psi$-field momentum density:
\begin{equation}
g_\psi = -\frac{1}{4\pi G}\,\mu\,\dot\psi\,\nabla\psi\,.
\end{equation}
This carries a factor of $1/G$, making it potentially large. However, this formula was derived from the
\emph{gravitational} field equation, where $\psi$ is sourced with $G/c^4$ coupling.

For the above-threshold mechanism, $\delta\psieff = \kappa(\eta - \etac)$ is \emph{not} a solution of the
gravitational field equation---it is a direct metric modification. The standard $\psi$-momentum formula
may not apply.

\subsubsection{Two sub-cases}

\textbf{Case A1: $\psi$-momentum formula applies.} Then $g_\psi \propto 1/G$, and the $\psi$-field can
carry large recoil momentum. The device accelerates one way; a $\psi$-wave propagates the other way
carrying the recoil. This is the ``$\psi$-drive'' concept from the interstellar travel analysis.
Net propulsion is possible.

\textbf{Case A2: $\psi$-momentum formula does not apply.} Then momentum conservation requires that the
EM stress tensor provides the recoil. But the EM stress tensor of a magnet is internal to the device.
The device cannot accelerate its own centre of mass using only internal EM stresses. Net propulsion
is \emph{not} possible in this case.

\subsection{Scenario (B): Violation of the equivalence principle}

If the EM source is somehow exempt from the $\psieff$ acceleration (e.g., because the modification is an
``interaction'' effect rather than a ``field'' effect), then:
\begin{itemize}
\item Test masses accelerate, but the magnet does not.
\item This violates the equivalence principle (WEP).
\item DFD explicitly assumes WEP, so this scenario is inconsistent with the theory.
\end{itemize}
We therefore rule out Scenario (B).

\subsection{Resolution: Momentum budget}

The resolution must come from a complete momentum budget. The key question is:

\emph{Does the above-threshold $\psieff$ modification carry its own momentum, independent of the
gravitational $\psi$-field momentum?}

If $\neff = \exp(\psieff)$ and $\psieff$ includes both the gravitational contribution and the threshold
contribution, then the \emph{total} $\psi$-field (including the threshold piece) should carry momentum.
The momentum density of the optical metric perturbation is:
\begin{equation}
g_{\text{optical}} \sim \frac{c^2}{8\pi G_{\text{eff}}}\,\partial_t(\delta\psieff)\,\nabla(\delta\psieff)\,,
\end{equation}
where $G_{\text{eff}}$ is the effective coupling for the threshold mechanism. If $G_{\text{eff}}$ is
\emph{not} Newton's $G$ but rather a coupling of order unity (consistent with $\kappa \sim 50$), then:
\begin{equation}
g_{\text{optical}} \sim c^2\,\partial_t(\delta\psieff)\,\nabla(\delta\psieff)\,,
\end{equation}
which can be large enough to carry the recoil momentum.

\textbf{Conclusion:} Propulsion is \emph{possible in principle} if:
\begin{enumerate}
\item The $\psieff$ field (including the threshold contribution) carries momentum.
\item The device creates an \emph{asymmetric} $\psieff$ profile (not symmetric about its centre).
\item $\psi$-waves or $\psieff$-perturbations propagate away from the device, carrying recoil momentum.
\end{enumerate}

This is consistent with the earlier ``$\psi$-bubble'' propulsion concept but requires the asymmetric
configuration of Scenario~2, not a symmetric field.

%=======================================================================
\section{Scenario 6: Practical Device Parameters}
\label{sec:scenario6}
%=======================================================================

\subsection{Threshold conditions}

From Eq.~\eqref{eq:rho_threshold}, the maximum gas density for threshold operation:

\begin{table}[H]
\centering
\begin{tabular}{rrrl}
\toprule
$B$ (T) & $\rho_{\mathrm{th}}$ (kg/m$^3$) & $P_{\mathrm{th}}$ (Pa) & $P_{\mathrm{th}}$ (Torr) \\
\midrule
10  & $2.43 \times 10^{-7}$ & $2.09 \times 10^{-2}$ & $1.57 \times 10^{-4}$ \\
20  & $9.71 \times 10^{-7}$ & $8.36 \times 10^{-2}$ & $6.27 \times 10^{-4}$ \\
30  & $2.18 \times 10^{-6}$ & $1.88 \times 10^{-1}$ & $1.41 \times 10^{-3}$ \\
50  & $6.07 \times 10^{-6}$ & $5.22 \times 10^{-1}$ & $3.92 \times 10^{-3}$ \\
\bottomrule
\end{tabular}
\caption{Threshold gas density and equivalent air pressure (at $T = 300$~K) for various magnetic field strengths.
These are moderate vacuum levels, achievable with standard rotary pumps.}
\label{tab:thresholds}
\end{table}

\subsection{10~T superconducting magnet at $\eta/\etac = 2$}

\begin{itemize}
\item $B = 10$ T (standard superconducting magnet)
\item $\rho = 1.21 \times 10^{-7}$ kg/m$^3$ (pressure $\sim 10^{-2}$~Pa, $\sim 10^{-4}$~Torr)
\item Bore radius $R = 0.1$ m
\end{itemize}

\begin{align}
\eta/\etac &= 2 \\
\delta\psieff &= 51.39 \times 1.824 \times 10^{-3} = 0.094 \\
|\nabla(\delta\psieff)| &\approx \frac{0.094}{0.1} = 0.94 \;\text{m}^{-1} \\
a &= \frac{c^2}{2} \times 0.94 = 4.2 \times 10^{16} \;\text{m/s}^2
\end{align}

\subsection{10~T at $\eta/\etac = 10$}

\begin{itemize}
\item $\rho = 2.43 \times 10^{-8}$ kg/m$^3$ ($P \sim 2 \times 10^{-3}$~Pa, $\sim 1.6 \times 10^{-5}$~Torr)
\end{itemize}

\begin{align}
\delta\psieff &= 0.844 \\
|\nabla(\delta\psieff)| &\approx 8.44 \;\text{m}^{-1} \\
a &= 3.79 \times 10^{17} \;\text{m/s}^2
\end{align}

This is the same $\delta\psieff$ as the solar corona calibration point, so it is within the
validated regime of the theory. The gradient length scale ($R = 0.1$~m) is much smaller than
the solar corona ($R_\odot \sim 7 \times 10^8$~m), which is why the acceleration is so much larger.

\subsection{50~T pulsed magnet}

\begin{itemize}
\item $B = 50$ T, $\rho = 10^{-7}$ kg/m$^3$, $R = 0.05$ m
\item $\eta/\etac = 60.7$, $\delta\psieff = 5.59$, $\exp(\delta\psieff) = 269$
\end{itemize}

\begin{align}
|\nabla(\delta\psieff)| &\approx 111.9 \;\text{m}^{-1} \\
a &= 5.0 \times 10^{18} \;\text{m/s}^2
\end{align}

\subsection{Energy in the bore}

For the 50~T scenario:
\begin{align}
V_{\text{bore}} &= \pi R^2 L = \pi \times (0.05)^2 \times 0.2 = 1.57 \times 10^{-3} \;\text{m}^3 \\
E_{\mathrm{EM}} &= U_{\mathrm{EM}} \times V_{\text{bore}} = 9.95 \times 10^8 \times 1.57 \times 10^{-3} = 1.56 \;\text{MJ} \\
m_{\text{gas}} &= \rho \times V_{\text{bore}} = 1.57 \times 10^{-10} \;\text{kg}
\end{align}

The gas mass is negligible. The stored EM energy is 1.56~MJ---substantial but not extraordinary.

\subsection{Summary: force vs.\ practicality}

\begin{table}[H]
\centering
\begin{tabular}{lccccc}
\toprule
\textbf{Configuration} & $B$ (T) & $\eta/\etac$ & $\delta\psieff$ & Regime & $a$ (m/s$^2$) \\
\midrule
10 T, $\eta/\etac = 2$   & 10 & 2    & 0.094 & Linear        & $4.2 \times 10^{16}$ \\
10 T, $\eta/\etac = 10$  & 10 & 10   & 0.844 & \textbf{Calibrated} & $3.8 \times 10^{17}$ \\
50 T pulsed               & 50 & 60.7 & 5.594 & Extrapolated  & $5.0 \times 10^{18}$ \\
Petawatt laser            & -- & $2 \times 10^5$ & 19{,}071 & Breakdown & $1.7 \times 10^{26}$ \\
\bottomrule
\end{tabular}
\caption{Summary of accelerations across scenarios. The ``Calibrated'' regime corresponds to the
solar corona anchor point. Note that even at $\delta\psieff = 0.094$ (just above threshold, firmly
in the linear regime), the acceleration is $10^{16}$~m/s$^2$.}
\label{tab:summary}
\end{table}

%=======================================================================
\section{Discussion: Why the Numbers Are So Large}
\label{sec:discussion}
%=======================================================================

\subsection{The $c^2$ enhancement}

The fundamental reason all numbers are enormous is the geodesic equation factor $c^2/2 = 4.49 \times 10^{16}$~m$^2$/s$^2$.
In standard gravity, this is cancelled by the weakness of gravitational potentials ($\psi \sim GM/(Rc^2) \sim 10^{-9}$
for laboratory masses). In DFD above threshold, $\delta\psieff$ is of order $\kappa\,\etac \sim 0.1$ even just above
threshold---\emph{thirteen orders of magnitude larger} than typical gravitational potentials.

The acceleration at marginal threshold ($\eta/\etac = 2$) over a 0.1~m gradient is:
\begin{equation}
a \sim \frac{c^2}{2} \times \frac{0.094}{0.1} \sim 4 \times 10^{16} \;\text{m/s}^2\,.
\end{equation}

\subsection{Comparison with known phenomena}

\begin{itemize}
\item Strongest continuous acceleration in laboratory: particle accelerators, $\sim 10^{15}$~m/s$^2$ (charged particles only)
\item Strongest gravitational field observed: neutron star surface, $\sim 10^{12}$~m/s$^2$
\item Predicted DFD threshold acceleration (marginal): $\sim 10^{16}$~m/s$^2$
\end{itemize}

This would be the strongest known force by a factor of $\sim 10^4$, applied to \emph{all matter universally}
(not just charged particles). It would have catastrophic effects on any matter in the fringe field region.

\subsection{Why this hasn't been seen}

There are two possible explanations:

\begin{enumerate}
\item \textbf{DFD is wrong} about the above-threshold mechanism. The light-bending prediction may be correct
      (it operates at $\delta\psieff < 1$) while the strong-field extrapolation fails.

\item \textbf{The vacuum condition is never met in existing experiments.} All high-field magnet experiments
      operate in air or helium at atmospheric pressure, where $\rho \sim 0.1$--$1.2$~kg/m$^3$.
      At these densities, even 50~T gives $\eta/\etac \sim 10^{-8}$---thirteen orders of magnitude
      below threshold. The threshold requires:
      \begin{itemize}
      \item 10~T: $P < 2 \times 10^{-2}$~Pa ($\sim 10^{-4}$~Torr)
      \item 50~T: $P < 0.5$~Pa ($\sim 4 \times 10^{-3}$~Torr)
      \end{itemize}
      No pulsed high-field magnet operates at these pressures, so the threshold has never been crossed
      in a controlled experiment.
\end{enumerate}

Explanation (2) makes an immediate prediction: \textbf{operate a pulsed magnet in vacuum}. If DFD is
correct, crossing the threshold would produce \emph{unmistakable} effects---anomalous forces on any
matter in the fringe field region.

\subsection{The hierarchy of predictions}

In order of theoretical confidence:
\begin{enumerate}
\item \textbf{Most confident:} Threshold \emph{exists}. Some anomaly occurs when $\eta > \etac$.
\item \textbf{Moderately confident:} The anomaly produces a $\psieff$ modification. Light-bending and
      clock rates change.
\item \textbf{Less confident:} The modification follows $\delta\psieff = \kappa(\eta - \etac)$ linearly
      for all $\eta > \etac$.
\item \textbf{Least confident:} The modification extends to $\eta/\etac \gg 10$ without saturation.
\end{enumerate}

\subsection{Recommended experimental strategy}

Given the uncertainty in the strong-coupling regime, the optimal experiment operates \emph{just above threshold}:
\begin{itemize}
\item $B = 10$~T (steady-state superconducting magnet)
\item $\rho \sim 10^{-7}$~kg/m$^3$ ($\eta/\etac \approx 2$)
\item $\delta\psieff = 0.094$---firmly in the linear regime
\item Look for anomalous light deflection (refractive index change $\sim 10\%$)
\item Look for anomalous clock rates (atomic clocks inside vs.\ outside the field region)
\item Look for anomalous forces on test masses at the fringe field boundary
\end{itemize}

Even $\delta\psieff = 0.094$ corresponds to a 10\% change in refractive index---easily measurable
with interferometry. This is the \emph{least} speculative prediction and should be tested first.

%=======================================================================
\section{Conclusions}
\label{sec:conclusions}
%=======================================================================

\begin{enumerate}
\item The DFD geodesic interpretation gives forces with coupling $\kappa \approx 51$ and no $G/c^4$
      suppression. The $c^2$ prefactor in the geodesic equation produces accelerations of
      $\sim 10^{16}$~m/s$^2$ even at marginal threshold.

\item The raw numbers are enormous---too large to ignore if correct, and too large to be plausible
      without experimental validation or a saturation mechanism.

\item The linear threshold formula is calibrated at $\eta/\etac \sim 10$ (solar corona). The 50~T
      pulsed magnet scenario ($\eta/\etac \sim 60$) is an extrapolation that may not hold.

\item The self-force problem for propulsion requires the $\psieff$ field (including the threshold
      contribution) to carry momentum. This is plausible but unproven.

\item The most robust prediction is the \emph{existence} of the threshold: a qualitative change in
      the optical metric when $\eta$ exceeds $\etac$ in a region of sufficient vacuum. This is
      testable with a 10~T superconducting magnet at $\sim 10^{-4}$~Torr.

\item No existing high-field magnet experiment operates in the required vacuum, which explains the
      absence of observed effects.

\item The first experiment should look for \emph{refractive index changes} (interferometry) or
      \emph{clock rate anomalies} (atomic clocks) at marginal threshold, not for propulsion forces.
      These are cleaner measurements and test the same underlying physics.
\end{enumerate}

%=======================================================================
\appendix
\section{Detailed Numerical Values}
\label{app:numbers}
%=======================================================================

\subsection{Physical constants used}
\begin{align}
c &= 2.998 \times 10^8 \;\text{m/s} \\
\mu_0 &= 4\pi \times 10^{-7} = 1.2566 \times 10^{-6} \;\text{T}\cdot\text{m/A} \\
\alpha &= 1/137.036 = 7.297 \times 10^{-3} \\
G &= 6.674 \times 10^{-11} \;\text{m}^3\,\text{kg}^{-1}\,\text{s}^{-2}
\end{align}

\subsection{DFD parameters}
\begin{align}
\kappa &= \frac{3}{8\alpha} = 51.39 \\
\etac &= \frac{\alpha}{4} = 1.824 \times 10^{-3} \\
\kappa\,\etac &= \frac{3}{8\alpha} \cdot \frac{\alpha}{4} = \frac{3}{32} = 0.09375
\end{align}

Note: $\kappa\,\etac = 3/32$ is an exact rational number. This means that at marginal threshold
($\eta = 2\etac$), $\delta\psieff = \kappa \cdot \etac = 3/32 = 0.09375$ exactly.

\subsection{Scenario 1: Complete computation}
\begin{align}
B &= 50 \;\text{T} \nonumber\\
U_{\mathrm{EM}} &= \frac{B^2}{2\mu_0} = \frac{(50)^2}{2 \times 1.2566 \times 10^{-6}} = 9.947 \times 10^8 \;\text{J/m}^3 \nonumber\\
\rho &= 10^{-7} \;\text{kg/m}^3 \nonumber\\
\rho c^2 &= 10^{-7} \times (2.998 \times 10^8)^2 = 8.988 \times 10^9 \;\text{J/m}^3 \nonumber\\
\eta &= \frac{9.947 \times 10^8}{8.988 \times 10^9} = 0.1107 \nonumber\\
\eta - \etac &= 0.1107 - 0.001824 = 0.1089 \nonumber\\
\delta\psieff &= 51.39 \times 0.1089 = 5.594 \nonumber\\
\exp(\delta\psieff) &= 268.7 \nonumber\\
R &= 0.05 \;\text{m} \nonumber\\
|\nabla\psieff| &\approx 5.594/0.05 = 111.9 \;\text{m}^{-1} \nonumber\\
a &= (c^2/2) \times 111.9 = 4.49 \times 10^{16} \times 111.9 = 5.03 \times 10^{18} \;\text{m/s}^2
\end{align}

\subsection{Threshold pressures for air at 300~K}

Using $P = \rho R_{\text{specific}} T$ with $R_{\text{specific}} = 287$~J/(kg$\cdot$K) for air:
\begin{equation}
P_{\text{th}} = \frac{B^2}{2\mu_0\,\etac\, c^2} \times 287 \times 300 = \frac{B^2 \times 8.61 \times 10^4}{2\mu_0\,\etac\, c^2}\,.
\end{equation}

\end{document}
